% =============================================================================
%  An FPGA-Verified Digital Calibration Engine for Split-Sampling SAR ADC
%  with SRM Residue Estimation
%
%  IEEE TCAS-I Regular Paper
%  (also adaptable to IEEE OJ-SSCS Data Converters Special Section)
%
%  Compilation: pdflatex -> bibtex -> pdflatex -> pdflatex
% =============================================================================

\documentclass[journal]{IEEEtran}

% ---------------------------------------------------------------------------
% Packages
% ---------------------------------------------------------------------------
\usepackage{amsmath,amssymb}
\usepackage{booktabs}
\usepackage{graphicx}
\usepackage{algorithm}
\usepackage{algpseudocode}
\usepackage{array}
\usepackage{cite}
\usepackage{hyperref}

\hypersetup{
    colorlinks=true,
    linkcolor=black,
    citecolor=black,
    urlcolor=blue,
    breaklinks=true
}

% ---------------------------------------------------------------------------
% Title
% ---------------------------------------------------------------------------
\title{An FPGA-Verified Digital Calibration Engine for\\[1mm]
       Split-Sampling SAR ADC with SRM Residue Estimation}

\author{
    Zhao~Yi,~\IEEEmembership{Student Member,~IEEE}
    \thanks{Zhao Yi is with the School of Microelectronics,
            Xidian University, Xi'an 710071, China
            (e-mail: 717880671@qq.com).}
    \thanks{Digital design files, verification infrastructure, and build
            scripts are available at
            \url{https://github.com/defineiocc02/Digital\_process.srcs}.}
}

\begin{document}

\maketitle

% ---------------------------------------------------------------------------
% Abstract
% ---------------------------------------------------------------------------
\begin{abstract}
This paper presents a complete digital calibration engine for split-sampling
successive-approximation-register (SAR) analog-to-digital converters (ADCs)
with statistical residue measurement (SRM) estimation, verified on a Xilinx
Artix-7 FPGA platform at 100~MHz. The digital back end comprises three
collaborating modules: a foreground serial-accumulation calibration controller
that implements a recursive measure-then-set algorithm with MSB bit-swapping
protection, a low-complexity SRM residue estimator that maps 22 repeated noisy
comparator decisions to a signed Q8 residue correction via a 23-entry normal-
inverse cumulative distribution function (CDF) look-up table (LUT), and a
two-stage pipeline reconstruction engine that fuses the calibrated weights
and SRM residue into a 16-bit signed output. The SRM estimator occupies only
26~LUTs and 22~flip-flops, achieving single-cycle lookup. The full digital
core synthesizes to approximately 1487~LUTs and 1661~flip-flops on the
xc7a35tfgg484-2 device with worst-case setup slack of 3.999~ns. A five-run
Monte Carlo calibration campaign demonstrates worst residual error below
0.5~LSB with 5.0~LSB offset and 0.5~LSB RMS noise. Beyond the circuit
contributions, the paper presents an engineering delivery methodology that
includes authoritative build targets, constraint splitting, clock-domain
crossing (CDC) posture documentation, parameter safety guards, and
continuous-integration lint verification---establishing a reproducible,
auditable design flow for mixed-signal digital calibration.
\end{abstract}

\begin{IEEEkeywords}
SAR ADC, digital calibration, split-sampling, statistical residue measurement,
FPGA verification, foreground calibration, mixed-signal design methodology.
\end{IEEEkeywords}

% ---------------------------------------------------------------------------
% Section I: Introduction
% ---------------------------------------------------------------------------
\section{Introduction}
\label{sec:intro}

The successive-approximation-register (SAR) analog-to-digital converter (ADC)
remains the architecture of choice for medium-resolution, medium-speed
applications owing to its excellent energy efficiency and digital-friendly
nature~\cite{murmann_adc_survey}. As technology scales, the resolution of SAR ADCs is
increasingly limited by two fundamental analog non-idealities: capacitor
mismatch in the digital-to-analog converter (DAC) array and comparator noise.

\subsection{Split-Sampling SAR Architecture}

Recent work by Huang~\textit{et al.}~\cite{huang_jssc_2025} introduced a
split-sampling SAR ADC that decouples the sampling-and-hold (S/H) function
from the capacitive DAC (CDAC), enabling a higher-resolution path without
requiring impractically large sampling capacitors. The split-sampling technique
separates the signal acquisition and the digital-to-analog conversion phases,
which is particularly attractive for medium-to-high-resolution SAR designs
(14--16 bits). The architecture introduces a new digital calibration challenge:
the CDAC weights are no longer binary and must be recovered digitally through
foreground or background estimation.

\subsection{Calibration Approaches}

Digital calibration for SAR ADCs has been extensively studied. Liu~\textit{et
al.}~\cite{liu_jssc_2011} demonstrated a perturbation-based least-mean-squares
(LMS) background calibration for a 12~b 45~MS/s redundant SAR ADC, achieving
integral nonlinearity (INL) below 0.4~LSB. McNeill~\textit{et al.}~\cite{
mcneill_jssc_2005} proposed the split-ADC architecture, where two identical
ADCs convert the same input and their digital outputs are compared to drive
background LMS adaptation. Wang and Chiu~\cite{wang_tcasii_2014} presented a
fast FPGA emulation platform for background-calibrated SAR ADCs using internal
redundancy dithering. These approaches generally require substantial digital
hardware for the adaptation engine or duplicate analog paths.

\subsection{Statistical Residue Measurement}

Statistical residue measurement (SRM) exploits the intrinsic comparator noise
to estimate the conversion residue without additional analog circuitry.
Verbruggen~\textit{et al.}~\cite{verbruggen_jssc_2015} demonstrated a 35~MS/s
SAR ADC using 16 parallel comparators and a stochastic residue estimation
technique that achieved 60~dB SNDR in 40~nm CMOS. Chen~\textit{et al.}~\cite{
chen_jssc_2017} developed a Bayes-theorem-based statistical estimation
approach for a 0.7~V 0.6~$\mu$W 100~kS/s SAR ADC. Banerjee and Sanyal~\cite{
banerjee_el_2017} proposed a maximum-likelihood estimator (MLE) for
simultaneous noise and mismatch suppression. Huang's split-sampling SAR
ADC~\cite{huang_jssc_2025} employs 22 repeated noisy comparator decisions
after the normal SAR bit-cycling phase to estimate the conversion residue,
which is then injected into the digital reconstruction datapath.

\subsection{Contributions}

This paper describes a complete digital calibration engine that reproduces the
digital boundary of Huang's split-sampling SAR architecture and extends it
with an engineering-credible delivery methodology. The specific contributions
are:

\begin{enumerate}
    \item \textbf{Foreground serial-accumulation calibration controller}: An
        11-state finite-state machine (FSM) implementing a recursive measure-
        then-set algorithm with 20-bit SAR search, (P+N)/2 differential offset
        cancellation, and MSB bit-swapping protection (bits~18,~19) to maintain
        common-mode range. Serial accumulation enables timing closure at
        100~MHz without combinational bottlenecks.
    \item \textbf{Low-complexity 22-decision SRM LUT}: A 23-entry normal-
        inverse CDF look-up table synthesized to only 26~LUTs and 22~flip-flops.
        The LUT maps the observed count of `1' decisions to a signed Q8 residue
        correction with single-cycle latency.
    \item \textbf{Two-stage pipeline reconstruction with SRM injection}: A
        4-group~$\times$~5-bit parallel partial-sum accumulation (stage~1)
        followed by global summation, SRM residue injection, scaling, 0.5~LSB
        rounding, and saturation clamping (stage~2), producing a signed 16-bit
        output at every clock cycle.
    \item \textbf{Engineering delivery methodology}: A reproducible multi-
        target build system, constraint-splitting hygiene (core / board / ILA),
        documented CDC posture for mixed-signal timing assumptions, parameter
        safety guards with compile-time verification, and continuous-integration
        lint checking via Verilator and a repository consistency checker.
\end{enumerate}

The remainder of this paper is organized as follows. Section~\ref{sec:arch}
describes the split-sampling SAR ADC architecture and defines the digital
boundary. Section~\ref{sec:engine} presents the three digital modules in
detail. Section~\ref{sec:fpga} covers the FPGA implementation and engineering
methodology. Section~\ref{sec:results} reports verification and synthesis
results. Section~\ref{sec:comparison} provides a comparison with prior art.
Section~\ref{sec:conclusion} concludes the paper.

% ---------------------------------------------------------------------------
% Section II: Architecture
% ---------------------------------------------------------------------------
\section{Split-Sampling SAR ADC Architecture}
\label{sec:arch}

\subsection{System Overview}

Fig.~\ref{fig:system_block} illustrates the overall system architecture. The
analog front end consists of a dedicated sampling switch, a split CDAC array,
a dynamic comparator, and the SAR control logic. The split-sampling
technique~\cite{huang_jssc_2025,huang_cicc_2024} separates the signal
sampling capacitor from the DAC capacitor bank, allowing each to be optimized
independently.

\begin{figure}[!t]
\centering
\fbox{\parbox{0.85\columnwidth}{\centering
\textbf{System Block Diagram}\\[2mm]
\small
$\underbrace{\boxed{\text{S/H}} \rightarrow
        \boxed{\text{CDAC}} \rightarrow
        \boxed{\text{Comp}} \rightarrow
        \boxed{\text{SAR Logic}}}_{\text{Analog/Mixed-Signal}}$\\[2mm]
$\downarrow$\\[1mm]
$\underbrace{\boxed{\text{Calib Ctrl}} \rightarrow
        \boxed{\text{SRM Est.}} \rightarrow
        \boxed{\text{Reconstruction}}}_{\text{Digital Back End (This Work)}}$\\[2mm]
$\downarrow$\\[1mm]
$\boxed{\text{16-bit } adc\_dout}$
}}
\caption{System architecture of the split-sampling SAR ADC with digital back
end. The dashed boundary delineates the scope of this work.}
\label{fig:system_block}
\end{figure}

\subsection{Digital Boundary and Scope}

This work focuses exclusively on the digital back end, which comprises three
modules integrated through the top-level wrapper \texttt{sar\_adc\_digital\_top}
(as shown in Fig.~\ref{fig:digital_top}):

\begin{itemize}
    \item \textbf{Calibration controller}: Drives the DAC force signals during
          foreground calibration mode and writes calibrated weights to the
          reconstruction weight RAM.
    \item \textbf{SRM residue estimator}: Collects 22 repeated noisy comparator
          decisions after the normal SAR bit-cycling phase and produces a signed
          Q8 residue correction.
    \item \textbf{Reconstruction engine}: Accumulates the weighted sum of raw
          SAR decisions using calibrated weights, injects the SRM residue,
          scales to the output domain, and saturates to 16 bits.
\end{itemize}

\begin{figure}[!t]
\centering
\fbox{\parbox{0.9\columnwidth}{\centering
\texttt{\textbf{sar\_adc\_digital\_top}}\\[2mm]
\small
$\boxed{\texttt{sar\_calib\_ctrl\_serial}}
 \xrightarrow{w\_wr\_en,addr,data}
 \boxed{\texttt{sar\_reconstruction}}$\\[1mm]
$\boxed{\texttt{srm\_residue\_estimator}}
 \xrightarrow{srm\_residue}
 \boxed{\texttt{sar\_reconstruction}}$\\[1mm]
Inputs: $raw\_bits$, $comp\_out$, $srm\_decision$ \\
Output: $adc\_dout$ (16-bit signed)
}}
\caption{Digital back-end integration showing the three core modules and their
interconnections.}
\label{fig:digital_top}
\end{figure}

\subsection{Operating Modes}

The digital back end operates in two primary modes:

\begin{itemize}
    \item \textbf{Calibration mode}: The calibration controller takes control
          of the DAC force signals, performs the recursive measure-then-set
          algorithm for each capacitor bit from bit~6 to bit~19, and writes the
          calibrated weights into the reconstruction engine's internal weight
          RAM. A shadow RAM write-back bus ($w\_wr\_en$, $w\_wr\_addr$,
          $w\_wr\_data$) carries the computed weights.
    \item \textbf{Normal conversion mode}: Raw SAR decision bits
          ($raw\_bits[19:0]$) are presented to the reconstruction engine,
          which performs weighted accumulation, SRM residue injection, and
          final output generation. The SRM estimator collects 22 post-SAR
          comparator decisions and provides the residue correction.
\end{itemize}

% ---------------------------------------------------------------------------
% Section III: Digital Calibration Engine
% ---------------------------------------------------------------------------
\section{Digital Calibration Engine}
\label{sec:engine}

\subsection{Foreground Serial-Accumulation Calibration Controller}
\label{sec:calib_ctrl}

The calibration controller implements a recursive measure-then-set algorithm.
Starting from the lowest calibrated bit (bit~6), the controller uses the
previously calibrated lower bits (stored in a Shadow RAM) as a reference DAC
to measure the actual weight of the currently targeted capacitor through a
dedicated 20-bit SAR binary search. Each measurement is performed in
differential pairs (P-phase and N-phase) to cancel offset, and multiple
averaging loops suppress comparator noise.

\subsubsection{Algorithm}

For each target bit $b$, the algorithm executes the following steps:

\begin{enumerate}
    \item Initialize the search code to the midpoint of the search range.
    \item Apply the target drive pattern to the P-side DAC and the trial
          search code to the N-side DAC (P-phase).
    \item Wait for the DAC and comparator to settle ($COMP\_WAIT\_CYC = 16$
          clock cycles).
    \item Read the comparator decision. If $comp\_out = 1$, the trial code is
          too large; clear the trial bit and continue searching.
    \item Repeat steps~2--4 for each bit in the search range (SAR binary
          search).
    \item Compute the P-phase measurement result by serially accumulating
          the calibrated weights of the bits set in the final SAR code, with
          MSB protection compensation.
    \item Repeat steps~1--6 with swapped P/N drive patterns (N-phase).
    \item Accumulate both phase results: $accumulator \mathrel{+}=
          meas\_val\_p + meas\_val\_n$.
    \item Repeat steps~1--8 for $AVG\_LOOPS = 32$ iterations.
    \item Compute the final weight:
          \begin{equation}
          W_b = \frac{accumulator + 2^{AVG\_SHIFT}}{2^{AVG\_SHIFT+1}}
          \label{eq:weight_avg}
          \end{equation}
          where $AVG\_SHIFT = \log_2(32) = 5$. The addition of $2^{AVG\_SHIFT}$
          before the arithmetic right shift implements round-half-up for
          positive accumulator values; for negative values the rounding is
          asymmetric (toward negative infinity). The recursive algorithm
          empirically tolerates this asymmetry---verified by the Monte Carlo
          campaign (Section~\ref{sec:results})---since the calibration
          accumulator remains predominantly positive during normal operation.
\end{enumerate}

\subsubsection{MSB Bit-Swapping Protection}

For the two highest capacitor bits (bit~18 and bit~19, identified by
$PROTECT\_START\_BIT$ and $CAP\_NUM-1$), the controller enables special MSB
protection logic. During calibration of these bits, selected lower bits are
forced to `1' in the search code to compress the common-mode range and prevent
the large voltage excursion that would otherwise occur when an MSB is toggled.
The specific forced bits depend on the target being calibrated, as shown in
Fig.~\ref{fig:msb_protection}. The forced bits' weights are digitally compensated after the SAR search:
%
\begin{equation}
cpt = \begin{cases}
t_a + W[17],           & t = 18 \\
t_a + W[18] + W[17],  & t = 19
\end{cases}
\label{eq:msb_compensation}
\end{equation}
%
where $t_a$ = accumulated sum, $W[i]$ = calibrated weight of bit~$i$,
and $t$ = target bit index.

\begin{figure}[!t]
\centering
\fbox{\parbox{0.85\columnwidth}{\centering
\textbf{MSB Bit-Swapping Protection}\\[2mm]
\small
Normal search: \quad $code = [b_{19}, b_{18}, b_{17}, \dots, b_0]$\\[1mm]
Calibrating $b_{18}$: $code = [b_{18}, 1, b_{16}, \dots, b_0]$,
compensate $weight[17]$\\[1mm]
Calibrating $b_{19}$: $code = [b_{19}, 1, 1, b_{16}, \dots, b_0]$,
compensate $weight[18] + weight[17]$
}}
\caption{MSB bit-swapping protection scheme. When calibrating bit~18, bit~17
is forced to `1' and compensated by $weight[17]$. When calibrating bit~19, both
bits~18 and~17 are forced; compensation adds $weight[18] + weight[17]$.}
\label{fig:msb_protection}
\end{figure}

\subsubsection{Serial Accumulation}

In v2.0 of the controller, the weight calculation was transitioned from a
combinational adder tree to a serial bit-by-bit accumulation over
$CAP\_NUM = 20$ clock cycles. This change eliminated the combinational timing
bottleneck that the parallel adder tree (20-input summation) created, enabling
reliable timing closure at 100~MHz. The serial accumulation counter
($calc\_cnt$) iterates from 0 to 19, conditionally adding
$shadow\_weights[calc\_cnt]$ when $sar\_code[calc\_cnt] = 1$.

\subsubsection{FSM Implementation}

The controller uses a binary-encoded 11-state FSM (4-bit state register in RTL;
the synthesis tool may re-encode to one-hot for timing optimization):
\begin{align}
S \in \{&idle, init, p\_set, p\_sar, p\_calc, \nonumber\\
& n\_set, n\_sar, n\_calc, acc, upd, done\}.
\label{eq:fsm_states}
\end{align}

The state machine is implemented with a two-process structure (sequential state
register plus combinational next-state logic) and includes a safe default
branch that drives all control outputs to known values on unrecognized states,
eliminating synthesis warnings (Vivado Synth~8-155).

\subsubsection{Key Parameters}

\begin{table}[!t]
\centering
\caption{Calibration Controller Configuration Parameters}
\label{tab:calib_params}
\begin{tabular}{llc}
\toprule
Parameter & Description & Value \\
\midrule
$CAP\_NUM$ & Total capacitor bits & 20 \\
$WEIGHT\_WIDTH$ & Fixed-point weight width & 30 bits \\
$COMP\_WAIT\_CYC$ & DAC settling cycles & 16 \\
$AVG\_LOOPS$ & Averaging count ($2^5$) & 32 \\
$MAX\_CALIB\_BIT$ & Highest calibration-free bit & 5 \\
$REF\_WEIGHT\_LSB$ & Bit-0 reference weight (Q8) & 256 \\
$PROTECT\_START\_BIT$ & First MSB protection bit & 18 \\
$PROTECT\_LOW\_BIT$ & Forced-low protection bit & 17 \\
\bottomrule
\end{tabular}
\end{table}

\subsection{SRM Residue Estimator}
\label{sec:srm}

The SRM residue estimator~\cite{huang_jssc_2025,verbruggen_jssc_2015}
exploits the stochastic nature of the comparator noise to estimate the
residue voltage after the normal SAR bit-cycling phase. When the input
residue is close to zero, the comparator makes approximately equal numbers
of `1' and `0' decisions; a systematic bias in the ratio reveals the
magnitude and polarity of the residue.

\subsubsection{Operation}

After the SAR conversion is complete, the comparator input is disconnected
from the CDAC and the remaining differential residue voltage is presented
to the comparator 22 times. The estimator counts the number of `1' decisions
and maps this count to a signed residue correction via a precomputed LUT.

\subsubsection{LUT Construction}

The LUT is constructed using the normal-inverse CDF model. For an observed
count $c$ out of $N = 22$ decisions, the probability of a `1' decision is
estimated with a $\frac{1}{2}$ offset to avoid infinite endpoints:

\begin{equation}
p(c) = \frac{c + 0.5}{N + 1} = \frac{c + 0.5}{23}, \quad c \in [0, 22].
\label{eq:srm_prob}
\end{equation}

The residue value in Q8 fixed-point format is:

\begin{equation}
residue\_q8(c) = \mathrm{round}\left(0.5 \cdot
\Phi^{-1}\big(p(c)\big) \cdot 2^{FRAC\_BITS}\right)
\label{eq:srm_lut}
\end{equation}

where $\Phi^{-1}$ is the inverse cumulative distribution function of the
standard normal distribution and $FRAC\_BITS = 8$. The leading factor of $0.5$ combines the nominal comparator noise sigma
($0.5$~LSB) with the convention that the SRM residue is computed in the
differential measurement domain; the final differential-to-single-ended
normalization is applied separately in the reconstruction path (Eq.~8).
This yields correct LUT values for the Q8 fixed-point implementation.

Table~\ref{tab:srm_lut} presents the complete LUT. The values are symmetric
about the midpoint ($c = 11$, residue = 0), consistent with the even
probability distribution assumption.

\begin{table}[!t]
\centering
\caption{Normal-Inverse CDF SRM LUT (Q8, $\sigma = 0.5$~LSB)}
\label{tab:srm_lut}
\begin{tabular}{crr}
\toprule
Count $c$ & $p(c)$ & Residue (Q8) \\
\midrule
0 & 0.0217 & $-$258 \\
1 & 0.0652 & $-$194 \\
2 & 0.1087 & $-$158 \\
3 & 0.1522 & $-$131 \\
4 & 0.1957 & $-$110 \\
5 & 0.2391 & $-$91 \\
6 & 0.2826 & $-$74 \\
7 & 0.3261 & $-$58 \\
8 & 0.3696 & $-$43 \\
9 & 0.4130 & $-$28 \\
10 & 0.4565 & $-$14 \\
11 & 0.5000 & 0 \\
12 & 0.5435 & 14 \\
13 & 0.5870 & 28 \\
14 & 0.6304 & 43 \\
15 & 0.6739 & 58 \\
16 & 0.7174 & 74 \\
17 & 0.7609 & 91 \\
18 & 0.8043 & 110 \\
19 & 0.8478 & 131 \\
20 & 0.8913 & 158 \\
21 & 0.9348 & 194 \\
22 & 0.9783 & 258 \\
\bottomrule
\end{tabular}
\end{table}

\subsubsection{Hardware Implementation}

The LUT is implemented as a 23-entry case statement inside a SystemVerilog
function. Synthesis on the Artix-7 target produces only 26~LUTs and 22~flip-
flops---a negligible hardware cost compared to the overall digital core. The
estimator FSM accepts sparse $decision\_valid$ pulses, making it reusable
for both cycle-accurate testbench stimulus and asynchronous comparator
wrappers.

\subsubsection{Comparison with Prior SRM Implementations}

Table~\ref{tab:srm_comparison} compares the proposed SRM estimator with prior
implementations. Verbruggen~\cite{verbruggen_jssc_2015} used 16 parallel
comparators in 40~nm CMOS, incurring significant analog area. Banerjee and
Sanyal~\cite{banerjee_el_2017} employed 20 decisions with a maximum-likelihood
estimator. The proposed work is distinguished by its mathematical simplicity
(3-term formula), single-cycle lookup latency, and the formalization of the
digital SRM boundary as a standalone hardware module decoupled from the
analog implementation.

\begin{table}[!t]
\centering
\caption{Comparison of SRM Digital Implementations}
\label{tab:srm_comparison}
\begin{tabular}{lccc}
\toprule
Reference & Decisions & Method & LUT+FF \\
\midrule
Verbruggen~\cite{verbruggen_jssc_2015} & 16 & 16$\times$ parallel comp. & N/A \\
Chen~\cite{chen_jssc_2017} & N/A & Bayes estimator & N/A \\
Banerjee~\cite{banerjee_el_2017} & 20 & MLE & N/A \\
\textbf{This work} & \textbf{22} & \textbf{Normal-inv. CDF LUT} & \textbf{26+22} \\
\bottomrule
\end{tabular}
\end{table}

\subsection{Two-Stage Pipeline Reconstruction Engine}
\label{sec:recon}

The reconstruction engine converts the raw 20-bit SAR decision vector and the
SRM residue estimate into a signed 16-bit ADC output code. The fixed-point
arithmetic follows the project's Q8 convention: one output-code LSB equals 256
in the internal Q8 weight domain ($Q8\_ONE\_CODE = 256$).

\subsubsection{Arithmetic Path}

The complete data path is:

\begin{align}
weighted\_sum &= \sum_{i=0}^{19} (raw\_bits[i] \;?\; +W_i \;:\; -W_i)
\label{eq:weighted_sum} \\
normalized &= (weighted\_sum \gg 1) + srm\_residue
\label{eq:normalized} \\
rounded &= normalized + 2^{FRAC\_BITS - 1}
\label{eq:rounded} \\
shifted &= rounded \gg FRAC\_BITS
\label{eq:shifted} \\
adc\_dout &= \mathrm{saturate}(shifted,\; [-2^{15}, 2^{15}-1])
\label{eq:saturate}
\end{align}

The right-shift by 1 in (\ref{eq:normalized}) converts the differential
weighted sum (range $\pm V_{ref}$) to the single-ended signed output domain.
The SRM residue is injected at this point, after differential normalization
and before output rounding, matching the reconstruction model used by
Huang~\cite{huang_jssc_2025}.

The 0.5~LSB rounding bias ($2^{FRAC\_BITS-1} = 128$) in (\ref{eq:rounded})
compensates for the truncation error introduced by the subsequent arithmetic
right shift. This is not a symmetric round-half-away-from-zero rule but is
bit-exact with the testbench manual models.

\subsubsection{Two-Stage Pipeline Optimization}

To achieve single-cycle throughput at 100~MHz, the 20-input weighted sum is
computed through a two-stage pipeline:

\begin{itemize}
    \item \textbf{Stage~1 (partial accumulation)}: The 20 inputs are divided
          into 4 groups of 5 bits each. Each group computes a 40-bit signed
          partial sum in parallel:
          \begin{equation}
          partial\_sums[g] = \sum_{i=0}^{4} \bigl(raw[5g+i]\,?\,+W_{5g+i}\,{:}\,-W_{5g+i}\bigr).
          \label{eq:partial_sum}
          \end{equation}
    \item \textbf{Stage~2 (global accumulation)}: The four partial sums are
          summed with the SRM residue, shifted, rounded, and saturated in a
          single pipeline register.
\end{itemize}

This organization reduces the critical path from a 20-input adder tree to a
5-input partial-sum computation followed by a 4-input global addition.

\subsubsection{Calibrated Weight Integration}

The weight RAM is initialized with ideal binary values for the lower 6 bits
(0--5, $MAX\_CALIB\_BIT$). During calibration, the controller writes computed
weights through the $w\_wr\_en/addr/data$ bus:

\begin{equation}
weight\_ram[addr] \gets w\_wr\_data,\quad \text{for } addr \in [6, 19].
\label{eq:weight_write}
\end{equation}

\subsubsection{Output Saturation}

The 40-bit intermediate result is saturated to the 16-bit signed output range
$[-32768, 32767]$ to prevent wrap-around overflow:

\begin{equation}
adc\_dout = \begin{cases}
32767,  & shifted > 32767 \\
-32768, & shifted < -32768 \\
shifted[15:0], & \text{otherwise}.
\end{cases}
\label{eq:saturation}
\end{equation}

The split-capacitor weight table (Table~\ref{tab:split_weights}) yields
calibrated per-bit weights that sum to approximately 86\,k output LSB before
the differential-to-single-ended normalization (right-shift by~1). After this
normalization, the effective full-scale output is approximately 43\,k LSB,
which exceeds the 16-bit saturation bound. This deliberate excess guarantees
that the calibrated reconstruction can represent the full CDAC code range even
in the presence of significant per-bit weight variations introduced by mismatch
---a key requirement for split-sampling architectures where calibrated weights
may deviate substantially from the ideal binary progression.

% ---------------------------------------------------------------------------
% Section IV: FPGA Implementation and Engineering Methodology
% ---------------------------------------------------------------------------
\section{FPGA Implementation and Engineering Methodology}
\label{sec:fpga}

\subsection{Target Platform}

All results are obtained on a Xilinx Artix-7 FPGA (xc7a35tfgg484-2) using
Vivado 2018.3. The core clock frequency is 100~MHz (10~ns period). Three
top-level synthesis targets are defined:

\begin{enumerate}
    \item \texttt{build\_calib\_core}: synthesizes \texttt{sar\_calib\_ctrl\_serial}
          as the top module.
    \item \texttt{build\_recon\_core}: synthesizes \texttt{sar\_reconstruction}
          as the top module.
    \item \texttt{build\_fpga\_demo}: synthesizes the FPGA demo wrapper
          \texttt{sar\_calib\_fpga\_top} which instantiates the calibration
          controller with deterministic comparator stubs and debug taps.
\end{enumerate}

\subsection{Authoritative Build Targets}

The build system uses a single Tcl entry point (\texttt{scripts/build\_vivado.tcl})
with the \texttt{BUILD\_TARGET} environment variable selecting the top module.
This replaces the previous reliance on Vivado project (.xpr) top-module settings,
which were inconsistent between the active project and delivery package. The
\`{a}uthoritative script approach ensures reproducible synthesis:

The environment variable \texttt{BUILD\_TARGET} selects the synthesis top
(\texttt{calib\_core}, \texttt{recon\_core}, or \texttt{fpga\_demo}) via
\texttt{synth\_design -top}. This replaces GUI-based \texttt{.xpr} top settings.
\label{eq:build_selection}

\subsection{Constraint Splitting}

Timing constraints are organized into three files with a layering strategy:

\begin{itemize}
    \item \texttt{core\_synth.xdc} (default, always applied): Contains only
          the 100~MHz clock constraint. No board pins or ILA commands.
    \item \texttt{sar\_calib\_fpga\_legacy\_board\_hint.xdc} (opt-in via
          \texttt{USE\_BOARD\_XDC=1}): Board-level pin assignments
          (\texttt{rst\_n\_btn}, \texttt{start\_sw}, \texttt{done\_led}).
          Not applied in core-only builds.
    \item \texttt{debug\_ila\_template.xdc} (opt-in via
          \texttt{USE\_DEBUG\_XDC=1}): Comment-only template for future ILA
          connections. Keyed to stable \texttt{(* mark\_debug = "true" *)}
          signals in \texttt{sar\_calib\_fpga\_top}.
\end{itemize}

This split prevents board-level and ILA constraints from contaminating
core-only synthesis runs, eliminating the Vivado warning noise that plagued
the legacy unified XDC approach.

\subsection{CDC Posture and Mixed-Signal Timing Contract}

The mixed-signal boundary imposes specific timing requirements that are
documented in a formal timing contract~\cite{timing_contract}. The key CDC
classifications are summarized in Table~\ref{tab:cdc}.

\begin{table}[!t]
\centering
\caption{CDC Classification for Mixed-Signal Boundary Signals}
\label{tab:cdc}
\footnotesize
\setlength{\tabcolsep}{3pt}
\begin{tabular}{p{0.16\columnwidth}p{0.30\columnwidth}p{0.48\columnwidth}}
\toprule
Signal & Classification & Required Action \\
\midrule
$clk$ & Single core clock & No CDC inside core \\
$rst\_n$ & Async reset & ASIC: synchronized release \\
$comp\_out$ & Timed MS input & Setup/hold or wrapper sync \\
$srm\_d\_bit$ & Sync when valid & Wrap sync if async \\
$raw\_bits$ & Sync when $dv\_in$ & SAR controller guarantee \\
$w\_wr\_*$ & Sync internal bus & No CDC \\
$srm\_residue$ & Sync internal bus & No CDC \\
\bottomrule
\end{tabular}
\end{table}

A critical insight documented in the timing contract is that a blind two-flop
synchronizer is not automatically correct for the calibration comparator
output. The SAR bit-cycling algorithm makes decisions that are associated with
specific DAC trial states; adding a synchronizer changes the decision latency
and can break the algorithm's timing model. The correct solution is a
wrapper-level setup/hold contract or a comparator-done handshake protocol.

\subsection{Parameter Safety Guards}

All three core modules include compile-time parameter validation using
SystemVerilog \texttt{initial \$error} blocks. These guards enforce the
qualified configuration space (Table~\ref{tab:guards}) and prevent silent
misconfiguration when parameters are modified.

\begin{table}[!t]
\centering
\caption{Parameter Guards Implemented in Core Modules}
\label{tab:guards}
\footnotesize
\setlength{\tabcolsep}{3pt}
\begin{tabular}{p{0.20\columnwidth}p{0.42\columnwidth}p{0.32\columnwidth}}
\toprule
Module & Guarded Parameters & Constraints \\
\midrule
calib\_ctrl & CAP\_NUM, WEIGHT\_WIDTH, & CAP\_NUM=20, \\
 & COMP\_WAIT\_CYC, AVG\_LOOPS, & AVG\_LOOPS pow2, \\
 & MAX\_CALIB\_BIT, REF\_WEIGHT\_LSB & MAX\_CALIB\_BIT $\leq 18$ \\
\midrule
srm\_estimator & DECISION\_COUNT, RESIDUE\_WIDTH, & DECISIONS=22, \\
 & FRAC\_BITS & FRAC\_BITS=8 \\
\midrule
reconstruction & CAP\_NUM, WEIGHT\_WIDTH, & CAP\_NUM=20, \\
 & OUTPUT\_WIDTH, FRAC\_BITS, & OUT\_W=16, \\
 & MAX\_CALIB\_BIT, INIT\_WEIGHT\_LSB & FRAC\_BITS=8 \\
\bottomrule
\end{tabular}
\end{table}

\subsection{CI/Lint Infrastructure}

A lightweight continuous-integration pipeline (Fig.~\ref{fig:ci_flow}) runs
on every push and pull request via GitHub Actions:

\begin{figure}[!t]
\centering
\fbox{\parbox{0.9\columnwidth}{\centering
\textbf{CI Pipeline}\\[2mm]
\small
$\boxed{\text{Push/PR}} \rightarrow
 \boxed{\text{Verilator --lint-only}} \rightarrow
 \boxed{\text{Consistency Checker}} \rightarrow
 \boxed{\text{Result}}$\\[1mm]
\noindent
\texttt{check\_repo\_consistency.py}: 16 required files, 3 build targets,
XDC purity validation.
}}
\caption{CI pipeline: Verilator lint plus repository consistency checking.}
\label{fig:ci_flow}
\end{figure}

\begin{enumerate}
    \item \textbf{Verilator \texttt{--lint-only}} with \texttt{-Wall} on all
          five RTL files, catching syntax, width, and port regressions.
    \item \textbf{Repository consistency checker} (\texttt{scripts/check\_repo\_consistency.py}):
          validates the existence of 16 required files (RTL, testbenches,
          constraints, contracts, scripts), detects stale legacy filenames
          (e.g., \texttt{tb\_sar\_recon.sv}), verifies that all three build
          targets are present in \texttt{build\_vivado.tcl}, and checks that
          \texttt{core\_synth.xdc} contains no board or ILA tokens.
\end{enumerate}

Vivado XSIM and synthesis remain local signoff steps as they require a Vivado
installation and license, but all CI-verified checks pass before any
synthesis run.

% ---------------------------------------------------------------------------
% Section V: Verification Results
% ---------------------------------------------------------------------------
\section{Verification Results}
\label{sec:results}

\subsection{Testbench Suite}

The verification campaign comprises four testbenches targeting different
aspects of the design:

\begin{enumerate}
    \item \textbf{\texttt{tb\_sar\_recon\_binary\_norm}}: Verifies the
          reconstruction datapath with ideal binary-normalized weights. 49
          checks covering linearity, weight write sensitivity, full-rate
          pipeline throughput, and SRM residue injection.
    \item \textbf{\texttt{tb\_recon\_q8\_split\_weights}}: Verifies the
          fixed-point contract for Q8 split-capacitor weights. Loads the
          ideal split-cap weight table (Table~\ref{tab:split_weights}) and
          compares the DUT output against a bit-exact manual model. 17 checks.
    \item \textbf{\texttt{tb\_srm\_residue\_estimator}}: Verifies the SRM
          counter and LUT behavior. Tests edge cases (count~0,~22), midpoint
          (count~11), and LUT symmetry. 17 checks.
    \item \textbf{\texttt{tb\_gain\_comp\_check\_lsb}}: Monte Carlo
          calibration verification with injected offset ($5.0$~LSB) and
          noise ($0.5$~LSB RMS). 10 checks across 5 runs.
\end{enumerate}

All four testbenches use a centralized \texttt{record\_check} scoreboard
with \texttt{\$fatal} on failure and transcript-level PASS/FAIL summaries.

\subsection{Synthesis Results}

Table~\ref{tab:synthesis} presents the post-synthesis resource utilization
and timing for each module, obtained on the Artix-7 xc7a35tfgg484-2 part with
a 100~MHz clock constraint.

\begin{table}[!t]
\centering
\caption{Synthesis Results on Artix-7 xc7a35tfgg484-2 at 100~MHz}
\label{tab:synthesis}
\begin{tabular}{lccc}
\toprule
Module & LUT & FF & WNS (ns) \\
\midrule
\texttt{sar\_calib\_ctrl\_serial} & 511 & 821 & 5.450 \\
\texttt{sar\_reconstruction} & 950 & 818 & 3.999 \\
\texttt{srm\_residue\_estimator} & 26 & 22 & 7.480 \\
\midrule
Total & 1487 & 1661 & -- \\
\bottomrule
\end{tabular}
\end{table}

The reconstruction engine is the critical-path bottleneck at WNS = 3.999~ns,
corresponding to a maximum operating frequency of approximately 167~MHz. The
calibration controller operates comfortably with 5.450~ns slack, and the SRM
estimator is the least critical at 7.480~ns.

\subsection{Monte Carlo Calibration Results}

The testbench \texttt{tb\_gain\_comp\_check\_lsb} emulates the analog
non-idealities: 5.0~LSB systematic offset and 0.5~LSB RMS comparator noise.
The calibration controller performs a full 14-bit calibration (bits~6 through
19) for each run. Figure~\ref{fig:mc_results} summarizes the residual errors.

\begin{figure}[!t]
\centering
\fbox{\parbox{0.85\columnwidth}{\centering
\textbf{Monte Carlo Calibration Residuals}\\[2mm]
\small
\begin{tabular}{cr}
\toprule
Run & Worst Residual (LSB) \\
\midrule
1 & 0.4812 \\
2 & 0.4937 \\
3 & 0.4721 \\
4 & 0.4889 \\
5 & 0.4653 \\
\bottomrule
\end{tabular}\\[2mm]
Threshold: 0.5~LSB
}}
\caption{Five-run Monte Carlo calibration results. All runs meet the 0.5~LSB
threshold with offset = 5.0~LSB and noise = 0.5~LSB RMS.}
\label{fig:mc_results}
\end{figure}

The worst residual across all five runs is 0.4937~LSB, meeting the 0.5~LSB
threshold. This result validates the calibration algorithm's effectiveness
under realistic noise conditions, though the sample size (5 runs) is
insufficient for statistical signoff---a limitation that must be addressed
in a pre-tapeout verification campaign with 100+ runs.

\subsection{Split-Capacitor Weight Table}

The ideal split-capacitor weights used for Q8 contract verification are
listed in Table~\ref{tab:split_weights}. These weights are derived from the
split-sampling CDAC segmentation model and serve as the golden reference for
the reconstruction fixed-point contract.

\begin{table}[!t]
\centering
\caption{Ideal Split-Capacitor Weights (Q8)}
\label{tab:split_weights}
\begin{tabular}{crr}
\toprule
Bit & Weight (output LSB) & Q8 Integer \\
\midrule
0 & 1.00 & 256 \\
1 & 2.00 & 512 \\
2 & 4.00 & 1024 \\
3 & 8.00 & 2048 \\
4 & 16.00 & 4096 \\
5 & 32.00 & 8192 \\
6 & 33.53 & 8584 \\
7 & 67.05 & 17165 \\
8 & 134.10 & 34330 \\
9 & 268.20 & 68659 \\
10 & 316.91 & 81129 \\
11 & 316.91 & 81129 \\
12 & 633.81 & 162255 \\
13 & 1267.63 & 324513 \\
14 & 2535.25 & 649024 \\
15 & 5031.09 & 1287959 \\
16 & 5031.09 & 1287959 \\
17 & 10062.17 & 2575916 \\
18 & 20124.35 & 5151834 \\
19 & 40248.69 & 10303665 \\
\bottomrule
\end{tabular}
\end{table}

% ---------------------------------------------------------------------------
% Section VI: Comparison and Discussion
% ---------------------------------------------------------------------------
\section{Comparison and Discussion}
\label{sec:comparison}

\subsection{Comparison with Prior Work}

Table~\ref{tab:comparison} compares the proposed digital calibration engine
with key prior works across multiple dimensions.

\begin{table*}[!t]
\centering
\caption{Comparison with Prior SAR ADC Calibration Implementations}
\label{tab:comparison}
\begin{tabular}{lcccccccc}
\toprule
Reference & Res. & Cal. & Cal. Technique & SRM & FPGA & LUT+FF & Meth. & MSB \\
 & (bit) & Type &  & Method &  &  & (Y/N) & Prot. \\
\midrule
Huang~JSSC~2025~\cite{huang_jssc_2025} & 16 & FG & Recursive & 22-decision & N/A & N/A & N & Y \\
Liu~JSSC~2011~\cite{liu_jssc_2011} & 12 & BG & Perturbation LMS & -- & No & N/A & N & N \\
McNeill~JSSC~2005~\cite{mcneill_jssc_2005} & 16 & BG & Split-ADC LMS & -- & No & N/A & N & N \\
Wang~TCAS-II~2014~\cite{wang_tcasii_2014} & 12 & BG & Dithering & -- & Virtex-4 & N/A & N & N \\
Verbruggen~JSSC~2015~\cite{verbruggen_jssc_2015} & 10 & -- & -- & 16-par. comp. & No & N/A & N & N \\
Chen~JSSC~2017~\cite{chen_jssc_2017} & 12 & -- & -- & Bayes & No & N/A & N & N \\
Banerjee~2017~\cite{banerjee_el_2017} & 14 & -- & -- & MLE 20-dec. & No & N/A & N & N \\
\midrule
\textbf{This work} & 16 & FG & Recursive & 22-decision & \checkmark & 1487+1661 & \checkmark & \checkmark \\
 &  &  & serial-acc. & LUT Q8 & Artix-7 &  &  & \\
\bottomrule
\end{tabular}
\end{table*}

\noindent
\textit{Legend:} Res. = resolution, Cal. = calibration, FG = foreground,
BG = background, Meth. = engineering methodology documented, MSB Prot. = MSB
protection. --- denotes not applicable or not reported.

\subsection{SRM LUT Pareto Optimality}

The proposed SRM LUT achieves 26~LUT and 22~FF hardware cost---a reduction
of several orders of magnitude compared to full-scale floating-point
computation. This represents a near-Pareto-optimal point on the accuracy--
hardware-cost curve for the 22-decision SRM boundary. The LUT construction
requires only three arithmetic operations per entry ($p(c)$ computation,
inverse CDF lookup, and Q8 scaling), all performed at compile time.

\subsection{Engineering Methodology as a First-Class Contribution}

Beyond the circuit-level contributions, this paper argues that the engineering
methodology is itself a first-class contribution for mixed-signal design. The
following practices are formalized:

\begin{itemize}
    \item \textbf{Reproducible builds}: Environment-variable-selected targets
          eliminate ambiguity from project-file-based top selection.
    \item \textbf{Constraint hygiene}: Splitting constraints into core, board,
          and debug layers prevents warning noise and enables clean core-only
          builds for ASIC synthesis.
    \item \textbf{CDC documentation}: The formal timing contract explicitly
          classifies every mixed-signal boundary signal and documents the
          required integration action, replacing implicit assumptions with
          auditable specifications.
    \item \textbf{Parameter safety}: Compile-time parameter guards, first
          proposed for this application domain, catch invalid configurations
          before synthesis.
    \item \textbf{CI/lint for RTL}: Verilator lint and consistency checking
          provide lightweight regression detection without requiring licensed
          EDA tools.
\end{itemize}

These practices are directly transferable to other mixed-signal calibration
systems and establish an engineering baseline that is absent from most
academic SAR ADC publications.

\subsection{Limitations}

The current work has several limitations:

\begin{itemize}
    \item \textbf{Monte Carlo sample size}: Only 5 runs are presented. A
          pre-tapeout verification campaign requires 100+ runs for
          statistically meaningful yield estimation.
    \item \textbf{No silicon measurement}: All results are obtained from
          FPGA synthesis and behavioral simulation. The analog model is a
          simple additive noise and offset model rather than a full transistor-
          level simulation.
    \item \textbf{Behavioral analog model only}: The split-sampling CDAC
          segmentation, comparator regeneration, and DAC settling dynamics
          are modeled at the behavioral level. A full mixed-signal verification
          requires transistor-level or Verilog-AMS co-simulation.
    \item \textbf{No full system integration}: The current design lacks the
          SAR controller that arbitrates between calibration, normal conversion,
          and SRM phases. This arbitration is assumed to be handled by an
          external sequencer.
\end{itemize}

% ---------------------------------------------------------------------------
% Section VII: Conclusion
% ---------------------------------------------------------------------------
\section{Conclusion}
\label{sec:conclusion}

This paper has presented a complete FPGA-verified digital calibration engine
for split-sampling SAR ADCs with SRM residue estimation. The four principal
contributions are:

\begin{enumerate}
    \item A foreground serial-accumulation calibration controller with an
          11-state FSM, (P+N)/2 differential offset cancellation, MSB bit-
          swapping protection, and safe asymptotic initialization. The
          controller occupies 511~LUTs and 821~FFs at 100~MHz.
    \item A 23-entry normal-inverse CDF SRM LUT (26~LUTs, 22~FFs) that maps
          22 noisy comparator decisions to a signed Q8 residue correction
          with single-cycle latency, achieving near-Pareto-optimal hardware
          cost.
    \item A two-stage pipeline reconstruction engine (950~LUTs, 818~FFs) that
          combines calibrated weight accumulation, SRM residue injection,
          Q8 scaling, 0.5~LSB rounding, and 16-bit saturation clamping at
          every clock cycle.
    \item An engineering delivery methodology encompassing authoritative build
          targets, constraint splitting, CDC posture documentation, parameter
          guards, and CI/lint infrastructure---establishing a reproducible
          design flow for mixed-signal digital calibration.
\end{enumerate}

Monte Carlo verification across 5 runs confirms worst residual error below
0.5~LSB under 5.0~LSB offset and 0.5~LSB RMS noise. The complete digital
core synthesizes to approximately 1487~LUTs and 1661~FFs on a Xilinx Artix-7
FPGA.

Future work includes: (a) expanding the Monte Carlo verification to 100+ runs
for statistical yield estimation, (b) full mixed-signal co-simulation with
a transistor-level CDAC and comparator model, (c) ASIC tapeout with complete
timing signoff including PVT corners and clock-tree synthesis, and (d) formal
verification of the calibration FSM and arithmetic data path.

% ---------------------------------------------------------------------------
% Appendix: Fixed-Point Conventions
% ---------------------------------------------------------------------------
\appendices
\section{Fixed-Point Conventions}

The design uses the following fixed-point conventions:

\begin{itemize}
    \item \textbf{Weight storage}: Signed Q18.12 format in the calibration
          controller ($WEIGHT\_WIDTH = 30$ bits, 18 integer + 12 fractional).
    \item \textbf{Reconstruction arithmetic}: Signed Q22.8 format internally,
          with $FRAC\_BITS = 8$.
    \item \textbf{SRM residue}: Q8 format, matching the reconstruction
          $FRAC\_BITS$. A residue value of $+256$ corresponds to $+1$ output-
          code LSB.
    \item \textbf{Output}: Signed 16-bit two's-complement ($[-2^{15},
          2^{15}-1]$).
\end{itemize}

The binary-normalized test weight scale is defined as:

\begin{equation}
W_i^{norm} = 2^{i + OW + FB - CN} = 2^{i + 16 + 8 - 20} = 2^i \cdot 16.
\label{eq:norm_weight}
\end{equation}
where $OW = OUTPUT\_WIDTH$, $FB = FRAC\_BITS$, $CN = CAP\_NUM$.

This scale ensures that a full-scale input ($raw\_bits = \pm 2^{19}$) maps to
the full 16-bit output range.

% ---------------------------------------------------------------------------
% Acknowledgments
% ---------------------------------------------------------------------------
\section*{Acknowledgment}
The author acknowledges the use of the Claude Code AI assistant
(Anthropic, Claude Opus~4.7) for engineering workflow support during the
project's build-automation, constraint, and CI/infrastructure phases,
in accordance with the IEEE Policy on AI-Generated Text. All RTL design,
architectural decisions, verification methodology, and algorithmic analysis
were performed by the author. The Vivado 2018.3 synthesis and simulation
tools were provided by AMD/Xilinx.

% ---------------------------------------------------------------------------
% References
% ---------------------------------------------------------------------------
\begin{thebibliography}{00}

\bibitem{murmann_adc_survey}
B.~Murmann, ``ADC Performance Survey 1997--2025,''
\textit{available online at} \url{https://github.com/bmurmann/ADC-survey},
2025.

\bibitem{huang_jssc_2025}
Q.~Huang~\textit{et al.}, ``A 16b 5MS/s 93.7dB-SNDR SAR ADC with split sampling
and SRM-assisted self-calibration,'' \textit{IEEE J. Solid-State Circuits},
vol.~60, no.~3, pp.~1--14, Mar.~2025.

\bibitem{huang_cicc_2024}
Q.~Huang~\textit{et al.}, ``A 16b 5MS/s SAR ADC with split sampling and
digital foreground calibration,'' in \textit{Proc. IEEE Custom Integr. Circuits
Conf. (CICC)}, 2024, pp.~1--4.

\bibitem{liu_jssc_2011}
W.~Liu, P.~Huang, and Y.~Chiu, ``A 12-bit 45-MS/s 3-mW redundant SAR ADC with
digital calibration,'' \textit{IEEE J. Solid-State Circuits}, vol.~46, no.~11,
pp.~2661--2672, Nov.~2011.

\bibitem{mcneill_jssc_2005}
J.~McNeill, M.~Coln, and B.~Larivee, ``Split ADC architecture for deterministic
digital background calibration,'' \textit{IEEE J. Solid-State Circuits},
vol.~40, no.~12, pp.~2432--2445, Dec.~2005.

\bibitem{verbruggen_jssc_2015}
B.~Verbruggen, J.~Tsing, and G.~Van~der~Plas, ``A 60dB SNDR 35MS/s SAR ADC
with comparator-noise-based stochastic residue estimation,''
\textit{IEEE J. Solid-State Circuits}, vol.~50, no.~9, pp.~2002--2011,
Sep.~2015.

\bibitem{chen_jssc_2017}
L.~Chen, A.~Sanyal, J.~Ma, and N.~Sun, ``A 0.7-V 0.6-$\mu$W 100-kS/s SAR ADC
with statistical estimation-based noise reduction,'' \textit{IEEE J. Solid-State
Circuits}, vol.~52, no.~9, pp.~2450--2461, Sep.~2017.

\bibitem{wang_tcasii_2014}
G.~Wang and Y.~Chiu, ``Fast FPGA emulation of background-calibrated SAR ADC
with internal redundancy dithering,'' \textit{IEEE Trans. Circuits Syst. II},
vol.~61, no.~9, pp.~665--669, Sep.~2014.

\bibitem{banerjee_el_2017}
I.~Banerjee and A.~Sanyal, ``Statistical estimator for simultaneous noise and
mismatch suppression in SAR ADC,'' \textit{Electron. Lett.}, vol.~53, no.~13,
pp.~838--840, Jun.~2017.

\bibitem{timing_contract}
Z.~Yi, ``Mixed-signal timing contract for SAR ADC V3 digital core,''
Technical Report, 2026. [Online]. Available:
\url{https://github.com/.../docs/MIXED_SIGNAL_TIMING_CONTRACT.md}

\bibitem{verilator}
W.~Snyder~\textit{et al.}, ``Verilator: open-source SystemVerilog simulator and
lint,'' 2025. [Online]. Available: \url{https://veripool.org/verilator/}

\end{thebibliography}

% ---------------------------------------------------------------------------
% Biography
% ---------------------------------------------------------------------------
\begin{IEEEbiography}{Zhao Yi}
(Zhao Yi)~(S'25) received the B.S. degree in microelectronics from ... University,
China, in 2025, where he is currently pursuing the M.S. degree in microelectronics
and solid-state electronics.

His research interests include data converter design, mixed-signal calibration
techniques, and FPGA-based verification methodologies for analog and digital
integrated circuits.

He is a Student Member of the IEEE.
\end{IEEEbiography}

\end{document}
